\chapter{IGF-1 (Insulin-like Growth Factor 1)}

\section*{What Is This Test?}
Insulin-like growth factor 1 (IGF-1, also called somatomedin C) is a 70-amino-acid single-chain polypeptide that mediates most of the growth-promoting and anabolic effects of growth hormone (GH). IGF-1 is produced primarily by the liver (75--90\% of circulating IGF-1) in response to GH binding to hepatic GH receptors, and also locally in bone, cartilage, muscle, and other peripheral tissues where it acts in an autocrine and paracrine manner. In the circulation, over 99\% of IGF-1 is bound to specific IGF-binding proteins (IGFBPs), primarily IGFBP-3 (which forms a 150 kDa ternary complex with IGF-1 and the acid-labile subunit, ALS), and to a lesser extent IGFBP-1, 2, 4, and 5. The ternary complex extends the half-life of IGF-1 from minutes (free IGF-1) to 15--20 hours (complexed). Because of this long half-life and the absence of pulsatile secretion (unlike GH), IGF-1 provides a stable, integrated measure of GH secretion over the preceding 24 hours. IGF-1 is the preferred screening test for acromegaly (GH excess) and is also used to screen for GH deficiency, monitor GH replacement therapy, and evaluate growth disorders. IGF-1 levels are strongly age-dependent: they are low at birth, rise during childhood and peak at puberty, and then decline progressively with age after the 3rd decade. For this reason, IGF-1 results MUST be interpreted using age- and sex-adjusted reference ranges (reported as absolute concentration and as a standard deviation score/Z-score relative to age-matched controls).

\section*{Alternative Names}
Somatomedin C; Insulin-like Growth Factor 1; IGF-I; Serum IGF-1; IGF-1, LC-MS/MS; IGF-1 Z-Score; IGF-1 SD Score

\section*{Principle}
IGF-1 is measured using competitive or two-site immunometric (chemiluminescent) immunoassays on automated analyzers. Most modern assays use a two-step format: first, IGF-1 is dissociated from its binding proteins using an acidic buffer containing excess IGF-2 (which displaces IGF-1 from IGFBPs but does not cross-react with the detection antibody). Then, a two-site sandwich immunoassay is performed with two monoclonal antibodies against different epitopes on IGF-1. One antibody (capture) is bound to a solid phase, and the second antibody (detection) is labeled with a chemiluminescent tag. The chemiluminescent signal is proportional to IGF-1 concentration. Common platforms: Roche Cobas ECLIA, Abbott Architect CMIA, Siemens IMMULITE/Atellica, and Diasorin Liaison. Inter-assay variability is substantial (CV 15--30\%) because of differences in dissociation protocols, antibody specificities, and calibration standards. In 2011, a consensus statement recommended standardization of IGF-1 assays to the recombinant WHO International Standard (IS 02/254). The gold standard reference method is liquid chromatography-tandem mass spectrometry (LC-MS/MS), which measures intact IGF-1 after trypsin digestion and quantification of signature proteotypic peptides. LC-MS/MS offers superior specificity (no cross-reactivity with IGF-2, pro-IGF-1, or IGFBPs) and is increasingly used in reference laboratories for IGF-1 measurement. The result should be reported as absolute concentration (ng/mL) AND as a Z-score (standard deviations from the age- and sex-matched mean), which is essential for clinical interpretation.

\section*{History}
IGF-1 was discovered in 1957 by William Daughaday and William Salmon at Washington University as "sulfation factor," a serum factor that mediated the effects of GH on cartilage sulfate incorporation. It was renamed "somatomedin" in the 1970s and subsequently "insulin-like growth factor 1" when its structural homology to proinsulin was recognized by Rinderknecht and Humbel in 1978. The first radioimmunoassay for IGF-1 was developed by Furlanetto and colleagues in 1977. The ternary complex (IGF-1/IGFBP-3/ALS) was characterized in the 1980s--1990s. The 2014 Endocrine Society guidelines established IGF-1 as the primary screening test for acromegaly, with the OGTT for GH as the confirmatory test. The 2011 consensus statement on GH and IGF-1 assay standardization addressed inter-assay variability and recommended Z-score reporting.

\section*{Why This Test Is Ordered}
\begin{itemize}
  \item Primary screening test for acromegaly in adults with clinical features (acral enlargement, coarsening of facial features, macroglossia, carpal tunnel syndrome, osteoarthritis, hypertension, diabetes, obstructive sleep apnea) --- elevated age-adjusted IGF-1 is the hallmark biochemical finding
  \item Monitoring disease activity after treatment of acromegaly (surgery, somatostatin analogues, pegvisomant, radiation) --- target IGF-1 is age-adjusted normal range
  \item Screening for GH deficiency in children with short stature (>2 SD below mean for age/sex, decreased growth velocity) and adults with known pituitary disease --- low IGF-1 supports GH deficiency but is not diagnostic alone
  \item Monitoring adequacy of GH replacement therapy in children and adults --- target IGF-1 in the age-adjusted mid-normal range (0 to +1 SD)
  \item Evaluating growth disorders including idiopathic short stature and IGF-1 deficiency (primary IGF-1 deficiency --- Laron syndrome or GH receptor mutations; secondary IGF-1 deficiency --- malnutrition, liver disease)
  \item Investigating nutritional status --- IGF-1 is sensitive to protein-calorie intake and declines rapidly with malnutrition, serving as a marker of nutritional adequacy
  \item Assessment of severe insulin resistance syndromes (IGF-1 may be low because insulin is needed for hepatic GH receptor expression)
  \item Serial monitoring in childhood cancer survivors who received cranial radiation --- acquired GH deficiency develops in 30--50\% and rising IGF-1 may indicate GH recovery
\end{itemize}

\section*{Is This a Primary or Secondary Test}
IGF-1 is the primary screening test for acromegaly (GH excess). A normal age-adjusted IGF-1 effectively excludes acromegaly (high negative predictive value). If IGF-1 is elevated, the OGTT with GH measurement is performed for confirmation. For GH deficiency, IGF-1 is a secondary test: a low IGF-1 supports the diagnosis but is not sufficient alone; GH stimulation testing is required for definitive diagnosis, because 30--40\% of adult GH-deficient patients have IGF-1 in the low-normal range. IGF-1 is also the primary monitoring marker for GH replacement therapy and acromegaly treatment.

\section*{How This Test Is Performed}
A healthcare professional draws a blood sample from a vein into a serum separator tube (gold-top) or plain red-top tube. Approximately 3--5 mL of blood is collected. The specimen is centrifuged, and serum is separated. IGF-1 is relatively stable: at room temperature for 24 hours, at 2--8 degrees C for 7 days, and frozen at -20 degrees C for 6 months. Because IGF-1 has minimal diurnal variation and is not pulsatile, the specimen can be drawn at any time of day and does not require fasting, though fasting is often part of a broader metabolic evaluation. The result must be reported as both absolute concentration (ng/mL) and age- and sex-adjusted Z-score (number of standard deviations from the mean for age and sex). A Z-score between -2 and +2 is considered normal.

\section*{What Preparation Is Needed}
No fasting is required specifically for IGF-1 measurement, but if the test is part of an evaluation that includes glucose or lipid profile, fasting may be needed. IGF-1 does not have significant diurnal variation, so timing of the draw is not critical. High-dose biotin (>5 mg/day) should be discontinued 48 hours before testing because of interference with streptavidin-biotin immunoassay platforms. Patients on GH therapy should have their IGF-1 measured 24 hours after their last GH injection (trough level). For patients on pegvisomant (GH receptor antagonist) for acromegaly, IGF-1 (not GH) is the marker used to titrate therapy, because pegvisomant blocks GH receptor signaling but GH levels rise paradoxically. Estrogen-containing oral contraceptives increase GH binding protein and decrease hepatic IGF-1 production (through first-pass effect on the liver), lowering IGF-1 levels by approximately 20--40\% in women taking oral estrogen. Transdermal estrogen has less effect on IGF-1. This must be considered when interpreting results: an oral contraceptive user with an IGF-1 at the lower limit of normal may have normal GH status.

\section*{Contraindications and Precautions}
The most important interpretive consideration is the strong age dependence of IGF-1. IGF-1 MUST be interpreted using age- and sex-adjusted reference ranges; an absolute concentration of 250 ng/mL is normal for a 20-year-old male but 3 SD above the mean for a 70-year-old male (indicating acromegaly). IGF-1 peaks during puberty (ages 13--17), with levels 2--3 times higher than adult values. After age 30, IGF-1 declines steadily by approximately 1--2\% per year. Factors decreasing IGF-1: malnutrition (starvation, anorexia nervosa, protein-calorie malnutrition), chronic liver disease (cirrhosis), uncontrolled diabetes mellitus (hepatic GH resistance due to low portal insulin), hypothyroidism (decreased hepatic GH receptor expression), chronic kidney disease, severe systemic illness, inflammatory bowel disease, and oral estrogen therapy. Factors increasing IGF-1: pregnancy (placental GH drives maternal IGF-1 production, IGF-1 increases 2--3 fold by third trimester), puberty (peak levels), hyperthyroidism, and obesity (mildly increased). IGF-1 may be normal in up to 30--40\% of adults with confirmed GH deficiency (false-normal), so a normal IGF-1 does not exclude adult GH deficiency. Conversely, low-normal or mildly low IGF-1 occurs in 5--10\% of the healthy population without GH deficiency. IGF-1 immunoassays have significant inter-assay variability (15--30\%). The same assay and laboratory must be used for serial monitoring. When changing platforms, a re-baseline IGF-1 should be measured before the transition. Macro-IGF-1 (IGF-1 bound to IgG antibodies, rare) can cause falsely elevated IGF-1 in some immunoassays.

\section*{Risks}
Minimal risk. Slight pain, bruising, or hematoma at venipuncture site. Rarely: vasovagal syncope, infection, or phlebitis. No radiation exposure. No risks from the test itself beyond standard venipuncture.

\section*{What Happens After the Test}
IGF-1 results are available within 1--4 hours on automated immunoassay platforms (or 3--7 days for LC-MS/MS send-out). Interpretation: IGF-1 Z-score >+2 (above 97.5th percentile for age and sex) is the hallmark of acromegaly. In a patient with clinical features, an elevated IGF-1 mandates an oral glucose tolerance test (75 g glucose) with GH measurements for confirmation. If the OGTT confirms acromegaly (GH nadir >1 ng/mL), pituitary MRI with gadolinium is the next step. If IGF-1 is normal (Z-score between -2 and +2), acromegaly is effectively excluded and further testing for GH excess is not needed. In GH deficiency evaluation: a low IGF-1 (Z-score <\-2) supports the diagnosis but is not sufficient alone; GH stimulation testing must be performed for definitive diagnosis. During GH replacement therapy: IGF-1 is measured every 3--6 months and the dose is adjusted to maintain IGF-1 in the mid-normal range (Z-score 0 to +1) for age and sex. During acromegaly treatment: IGF-1 is measured every 3--6 months; normalization of age-adjusted IGF-1 is the primary biochemical goal of treatment, correlating with reversal of soft tissue changes and reduced mortality.

\section*{Understanding Results}

\subsection*{Below Normal}
\begin{itemize}
  \item \textbf{Low IGF-1 (Z-score <\-2) in a child with short stature:} Suggests GH deficiency or GH insensitivity. Requires GH stimulation testing for confirmation. If GH peak is low (<10 ng/mL on 2 stimulation tests), GHD is confirmed and recombinant GH therapy is indicated. If GH peak is normal or elevated (>10 ng/mL), GH insensitivity (Laron syndrome, GH receptor mutations causing primary IGF-1 deficiency) should be considered, and recombinant IGF-1 (mecasermin) therapy may be indicated.
  \item \textbf{Low IGF-1 (Z-score <\-2) in an adult with pituitary disease:} Suggests adult GH deficiency. Approximately 30--40\% of adult GHD patients have IGF-1 in the low-normal range, so a normal IGF-1 does NOT exclude GHD. GH stimulation testing (ITT or GHRH-arginine) is required for definitive diagnosis. Adults with confirmed severe GHD benefit from GH replacement (improved body composition, bone density, quality of life, and cardiovascular risk profile).
  \item \textbf{Malnutrition/starvation:} IGF-1 is exquisitely sensitive to nutritional status. A 7--10 day fast reduces IGF-1 by 60--80\%, despite elevated GH (nutritional GH resistance). IGF-1 normalizes with refeeding. Used as a marker of nutritional adequacy in anorexia nervosa, inflammatory bowel disease, and critically ill patients.
  \item \textbf{Chronic liver disease (cirrhosis):} Decreased hepatic GH receptor expression and decreased IGF-1 synthesis. Low IGF-1 with elevated GH. IGF-1 levels correlate with severity of liver disease (Child-Pugh score).
  \item \textbf{Hypothyroidism:} Thyroid hormone is necessary for normal GH receptor expression in liver. Untreated hypothyroidism decreases GH secretion and hepatic IGF-1 production. IGF-1 normalizes with levothyroxine therapy.
  \item \textbf{Uncontrolled diabetes mellitus (Type 1, severe Type 2 with insulin deficiency):} Low portal insulin impairs hepatic GH receptor expression and decreases IGF-1 production despite elevated GH (hepatic GH resistance). IGF-1 normalizes with improved glycemic control.
  \item \textbf{Primary IGF-1 deficiency (Laron syndrome):} Autosomal recessive GH receptor mutations causing GH insensitivity. Elevated GH, low IGF-1, low IGFBP-3. Severe short stature, mid-face hypoplasia, obesity. Treatment with recombinant IGF-1 (mecasermin).
  \item \textbf{Oral estrogen therapy:} First-pass hepatic effect of oral estrogen decreases IGF-1 by 20--40\%. Transdermal estrogen has minimal effect.
\end{itemize}

\subsection*{Above Normal}
\begin{itemize}
  \item \textbf{Elevated IGF-1 (Z-score >+2) with clinical acromegaly features:} The hallmark biochemical finding in acromegaly. Requires confirmatory OGTT with GH measurement. IGF-1 integrates GH secretion and correlates with disease severity and soft tissue manifestations. IGF-1 >1.5--2 times the upper limit of normal is strongly indicative.
  \item \textbf{Acromegaly (GH-secreting pituitary adenoma):} IGF-1 elevated above age-adjusted upper limit; GH fails to suppress below 1 ng/mL after 75 g glucose. The degree of IGF-1 elevation correlates with tumor size and disease activity. IGF-1 normalizes after successful surgical resection or with somatostatin analogue (octreotide, lanreotide) therapy.
  \item \textbf{Pregnancy (second and third trimester):} Placental GH (GH-V) replaces pituitary GH and drives maternal IGF-1 production. IGF-1 rises 2--3 fold above non-pregnant values, peaking in the third trimester. This is physiologic and returns to normal after delivery.
  \item \textbf{Puberty (physiologic):} IGF-1 levels peak during mid-late puberty (Tanner stages III--V), when values are 2--3 times higher than adult normal ranges due to the synergistic effect of sex steroids and GH on IGF-1 production. Pubertal staging is essential for interpreting IGF-1 in adolescents.
  \item \textbf{Drug interference:} Recombinant human GH therapy --- IGF-1 increases dose-dependently; target is mid-normal range. Exogenous IGF-1 administration (mecasermin for severe primary IGF-1 deficiency) directly elevates measured IGF-1. Insulin (endogenous or exogenous) can modestly increase IGF-1 through upregulation of hepatic GH receptor.
  \item \textbf{Growth hormone excess without acromegaly (gigantism):} Markedly elevated IGF-1 in children/adolescents with GH excess before epiphyseal closure, causing excessive linear growth. Causes: pituitary gigantism (GH-secreting adenoma), McCune-Albright syndrome, Carney complex.
  \item \textbf{Excess GH therapy (overtreatment):} IGF-1 above the age-adjusted normal range indicates excessive GH replacement dose. GH dose should be decreased. Chronically elevated IGF-1 in GH-treated patients may confer acromegaly-like risks (fluid retention, joint pain, insulin resistance, theoretical increased neoplasia risk).
  \item \textbf{Hyperthyroidism:} Thyroid hormone excess increases hepatic GH receptor expression and IGF-1 production. IGF-1 can be mildly elevated (Z-score +1 to +2) in Graves disease or toxic nodular goiter. Normalizes with treatment.
\end{itemize}

\section*{Normal Results}
\begin{itemize}
  \item \textbf{IGF-1 (Males, age 18--19):} 170--520 ng/mL (Z-score -2 to +2)
  \item \textbf{IGF-1 (Females, age 18--19):} 150--460 ng/mL
  \item \textbf{IGF-1 (Males, age 20--24):} 130--420 ng/mL
  \item \textbf{IGF-1 (Females, age 20--24):} 120--380 ng/mL
  \item \textbf{IGF-1 (Males, age 25--29):} 110--370 ng/mL
  \item \textbf{IGF-1 (Females, age 25--29):} 100--350 ng/mL
  \item \textbf{IGF-1 (Males, age 30--39):} 90--310 ng/mL
  \item \textbf{IGF-1 (Females, age 30--39):} 85--290 ng/mL
  \item \textbf{IGF-1 (Males, age 40--49):} 80--270 ng/mL
  \item \textbf{IGF-1 (Females, age 40--49):} 75--250 ng/mL
  \item \textbf{IGF-1 (Males, age 50--59):} 70--240 ng/mL
  \item \textbf{IGF-1 (Females, age 50--59):} 65--230 ng/mL
  \item \textbf{IGF-1 (Males, age 60--69):} 60--210 ng/mL
  \item \textbf{IGF-1 (Females, age 60--69):} 55--200 ng/mL
  \item \textbf{IGF-1 (Males/Females, age >70):} 50--180 ng/mL
  \item \textbf{IGF-1 (Children, pubertal peak):} 200--900 ng/mL (2--3 times higher than adult, peaking at Tanner III--IV)
  \item \textbf{IGF-1 (Newborns):} 20--100 ng/mL (low at birth, rising through infancy)
  \item \textbf{Note:} These approximate ranges are for illustrative purposes. Actual reference ranges are assay- and laboratory-specific and MUST be used for clinical interpretation. Z-score transformation normalizes for age and sex and should always be provided with the result.
\end{itemize}

\section*{Follow-up Tests}
\begin{itemize}
  \item \textbf{Oral glucose tolerance test with GH (OGTT-GH):} Gold standard for confirming acromegaly when IGF-1 is elevated. 75 g glucose orally; GH measured at 0, 30, 60, 90, 120 minutes. Failure of GH to suppress below <1 ng/mL (<0.4 ng/mL with ultrasensitive assay) confirms acromegaly.
  \item \textbf{Pituitary MRI with gadolinium contrast:} Localizes GH-secreting pituitary adenoma in confirmed acromegaly. Thin-section (1--3 mm) coronal and sagittal T1-weighted images before and after contrast. Macroadenoma (>1 cm) found in 75--80\%, microadenoma in remainder.
  \item \textbf{GH stimulation testing (ITT, GHRH-arginine, glucagon):} To confirm GH deficiency when IGF-1 is low or clinical suspicion is high despite normal IGF-1. Gold standard is ITT.
  \item \textbf{IGFBP-3 (IGF-binding protein 3):} Main carrier protein for IGF-1 in circulation (ternary complex with ALS). Also GH-dependent but less sensitive than IGF-1. Low IGFBP-3 with low IGF-1 supports GH deficiency; normal IGFBP-3 with low IGF-1 suggests nutritional cause.
  \item \textbf{ALS (acid-labile subunit):} Also GH-dependent. Used for research and complex pediatric cases.
  \item \textbf{GH receptor gene (GHR) sequencing:} If primary IGF-1 deficiency (Laron syndrome) is suspected (low IGF-1, elevated or normal GH, poor response to GH stimulation). Confirms GH receptor mutations.
  \item \textbf{Bone age X-ray (left hand/wrist):} In children with short stature and low IGF-1. Delayed bone age is characteristic of GH deficiency. Advanced bone age occurs in precocious puberty.
  \item \textbf{Colonoscopy:} For patients with confirmed acromegaly at diagnosis, because of 2--3 fold increased risk of colorectal adenomas/carcinoma.
  \item \textbf{Echocardiogram:} Assesses acromegalic cardiomyopathy (biventricular hypertrophy, diastolic dysfunction).
\end{itemize}

\section*{References}
\begin{enumerate}
  \item Katznelson L, Laws ER Jr, Melmed S, et al. Acromegaly: an Endocrine Society clinical practice guideline. J Clin Endocrinol Metab. 2014;99(11):3933-3951.
  \item Melmed S. Acromegaly. N Engl J Med. 2019;380(15):1462-1471.
  \item Grimberg A, DiVall SA, Polizzotto C, et al. Guidelines for growth hormone and insulin-like growth factor-I treatment in children and adolescents. Horm Res Paediatr. 2016;86(6):361-397.
  \item Molitch ME, Clemmons DR, Malozowski S, et al. Evaluation and treatment of adult growth hormone deficiency: an Endocrine Society clinical practice guideline. J Clin Endocrinol Metab. 2011;96(6):1587-1609.
  \item Clemmons DR. Consensus statement on the standardization and evaluation of growth hormone and insulin-like growth factor assays. Clin Chem. 2011;57(4):555-559.
  \item Le Roith D. Seminars in medicine of the Beth Israel Deaconess Medical Center: insulin-like growth factors. N Engl J Med. 1997;336(9):633-640.
  \item Salmon WD Jr, Daughaday WH. A hormonally controlled serum factor which stimulates sulfate incorporation by cartilage in vitro. J Lab Clin Med. 1957;49(6):825-836.
  \item Burtis CA, Ashwood ER, Bruns DE. Tietz Textbook of Clinical Chemistry and Molecular Diagnostics. 6th ed. Elsevier; 2023.
\end{enumerate}

\clearpage
